\subsection{A comparison in degree five}
For a genus-two core, the first prime beyond the cubic gives
\begin{align*}
 D_{2,5}(w)=w^6\big(&444088+142232w+340864w^2-156542w^3+75396w^4\\
                    &-23737w^5+5054w^6-696w^7+56w^8-2w^9\big).
\end{align*}
Here the subscripts denote genus and covering degree as in \eqref{eq:compact-prime}. For example, at $5\alpha$ in output genus
seven the scalar exponent is the elementary exponent minus $444088$
times the sum of the genus-two elementary exponents at its effective
fifth roots. Larger core genera are excluded by Theorem~\ref{thm:main}.

For the calculation let $Z_a$ denote the disconnected degree-$a$
residue and $L_a$ the coefficient of its logarithm. Apply
$aL_a=aZ_a-\sum_{b=1}^{a-1}b L_b Z_{a-b}$ at each pair of weights.
For a genus-two core $Z_a$ is the trace of the balanced handle
of \cite[\S\S5--8]{BP}. In degrees $1$ through $5$ the transposition
operator is cyclic on the unit over $\QQ(r)$, $r=\exp(iu)$, at both
weight specializations. Using the powers of that operator on the unit
as basis, express each partition class as a polynomial in the operator.
Multiply the handle and the inverse level operators modulo its
characteristic polynomial and take the trace. At antidiagonal weights
one may equivalently sum $\prod_{b\in\lambda}(2-r^{\ell(b)}-
r^{-\ell(b)})$ over partitions, with $\ell(b)$ the hook length.
Subtract the two values of $L_5$ and substitute $w=2-r-r^{-1}$.
These are rational function identities over $\QQ$.
